\section{Conclusions and Future Work}\label{sec:conclusions}

The paper designed a cluster selection strategy that adds client-to-cluster latency to Karmada's scheduling without replacing Karmada's own stages. Regarding RQ1, the analysis of the Karmada v1.19.1 source code showed that a custom objective has to be computed inside a single score plugin, that it only takes effect together with a cluster spread constraint, and that Karmada's default behaviour for a new single-cluster workload already prefers the cluster with the most available replicas. Regarding RQ2, mandatory conditions are evaluated in the filter stage, and the capacity condition reuses Karmada's replica estimators so that it agrees with the selection stage. Regarding RQ3, a network measurement contributes to the score only if it is recent and based on enough samples; a missing measurement is treated as the least favourable value, and the score falls back to resource-only ranking when no candidate has a valid measurement. Regarding RQ4, the design trades global optimality and joint node selection for compatibility with Karmada and a closed-form score whose decisions can be predicted, and it inherits the limitations of request-based capacity information.

Future work consists of implementing the plugins and the latency probe, and carrying out the experimental validation of \cref{sec:validation} against the default and resource-only configurations. Further extensions are the rebalancing of workloads when conditions change, which Karmada supports only through explicit triggers \cite{karmada_docs_rebalancer}; the use of measured utilisation alongside requested resources; several client populations; latency between services for workloads split across clusters; and an evaluation of performance, scalability and reliability on larger installations.
